\section{Planteamiento del problema}


% Aspectos generales del estudio
% Problemática a niverl mundial
% Problemática a nivel regional
% Problemática a nivel nacional
% Problemática a nivel estatal
% Problemática a nivel local y/o objeto de estudio

La movilidad independiente es un elemento fundamental para la autonomía y calidad de vida de las personas. Sin embargo, para las personas con discapacidad visual o visión reducida, desplazarse de manera segura en entornos urbanos representa un desafío constante debido a la dificultad para percibir obstáculos, identificar rutas adecuadas y orientarse en espacios desconocidos. De acuerdo con la Organización Mundial de la Salud, más de 2,200 millones de personas en el mundo presentan algún grado de discapacidad visual, lo que evidencia la magnitud de esta problemática a nivel global (OMS, 2023). \cite{OMS_ceguera}

\\

En el caso de las personas con visión reducida o ceguera, la movilidad cotidiana suele depender de herramientas tradicionales como el bastón blanco, el perro guía o la asistencia de otras personas. Aunque el bastón es una herramienta ampliamente utilizada, su alcance se limita a detectar obstáculos cercanos al usuario y no proporciona información anticipada sobre el entorno ni orientación en rutas específicas. Esto puede generar situaciones de riesgo al transitar por calles, cruzar avenidas o desplazarse en espacios con múltiples obstáculos.

\\

En los últimos años se han desarrollado diversas tecnologías orientadas a mejorar la accesibilidad para personas con discapacidad visual, como aplicaciones móviles de navegación, sistemas de reconocimiento de voz y herramientas de lectura de texto mediante cámaras.

\\

Paralelamente, el avance de la inteligencia artificial ha permitido el desarrollo de sistemas capaces de interpretar información visual mediante técnicas de visión por computadora y aprendizaje profundo. Modelos de detección de objetos basados en redes neuronales convolucionales, como YOLO (You Only Look Once), han demostrado ser capaces de identificar múltiples objetos dentro de una imagen en tiempo real con altos niveles de precisión. Estas tecnologías han sido aplicadas en áreas como la conducción autónoma, la robótica y la vigilancia inteligente, lo que demuestra su potencial para analizar entornos complejos.

\\

Asimismo, el desarrollo de dispositivos de bajo costo y alto rendimiento, como la Raspberry Pi, ha facilitado la implementación de sistemas inteligentes portátiles capaces de procesar información en tiempo real mediante cámaras y sensores. La integración de estos dispositivos con tecnologías de reconocimiento de voz y de mapeo del entorno abre la posibilidad de desarrollar soluciones accesibles que permitan a los usuarios interactuar con el sistema mediante comandos de voz y recibir indicaciones auditivas sobre su entorno.

\\

A pesar de estos avances tecnológicos, todavía existe una brecha significativa entre las capacidades actuales de la inteligencia artificial y su aplicación práctica en sistemas de asistencia para personas con discapacidad visual. Muchas de las soluciones existentes no integran de manera eficiente la detección de obstáculos, la orientación en el entorno y la interacción auditiva dentro de un mismo sistema portátil que pueda ser utilizado en la vida cotidiana.

\\

En este contexto surge la necesidad de desarrollar un sistema que combine tecnologías de visión artificial, inteligencia artificial, sensores y mapeo del entorno para asistir a personas con visión reducida durante su desplazamiento. Un sistema de este tipo podría capturar imágenes del entorno mediante una cámara, analizar la información visual utilizando algoritmos de detección de objetos, identificar obstáculos relevantes y proporcionar indicaciones auditivas al usuario sobre la presencia de objetos o rutas a seguir.

\\

Además, la incorporación de tecnologías de reconocimiento de voz permitiría que el usuario interactúe con el sistema de manera natural, solicitando indicaciones de navegación o información sobre su entorno sin necesidad de utilizar interfaces visuales. La integración de estos elementos en un prototipo basado en hardware portátil podría contribuir a mejorar la seguridad, autonomía y calidad de vida de las personas con discapacidad visual.

\\

Por lo tanto, el problema central que aborda esta investigación se relaciona con la falta de sistemas accesibles e inteligentes que permitan a las personas con visión reducida desplazarse de manera más segura y autónoma mediante el uso de tecnologías de inteligencia artificial, visión por computadora y navegación auditiva en tiempo real.

\\

Ante esta problemática, el presente trabajo propone el desarrollo de un sistema de navegación asistida basado en inteligencia artificial que permita detectar obstáculos mediante visión artificial, procesar la información en tiempo real y proporcionar indicaciones auditivas al usuario para facilitar su desplazamiento en diferentes entornos.

\subsection{Preguntas de investigación}
Con base en la problemática identificada sobre las dificultades de movilidad y orientación que enfrentan las personas con discapacidad visual, surgen diversas interrogantes que orientan el desarrollo de la presente investigación. A nivel mundial, millones de personas presentan algún grado de pérdida de visión, y en México se estima que alrededor de 11 millones de personas viven con ceguera o discapacidad visual, lo que refleja la relevancia social de desarrollar tecnologías de apoyo que faciliten su desplazamiento y autonomía.\\
Asimismo, en el país se han registrado aproximadamente 8.8 millones de personas con algún tipo de discapacidad, lo que representa cerca del $7.2\%$ de la población, evidenciando la necesidad de implementar soluciones tecnológicas inclusivas que contribuyan a mejorar la calidad de vida de este sector.
\begin{itemize}
    \item ¿Qué es la inteligencia artificial?
    \item ¿Qué es la vista?
    \item ¿Qué es discapacidad visual?
    \item ¿Qué es un sistema de navegación?
    \item ¿Qué es persona con discapacidad visual?
    \item ¿Qué es asistencia?
    \item ¿Qué es la localización y mapeo simultáneo (SLAM)?
    \item ¿Qué es ruta?
    \item ¿De qué manera puede desarrollarse un sistema de navegación asistida que utilice inteligencia artificial, visión por computadora y guía auditiva para mejorar la movilidad de personas con visión reducida en entornos urbanos?
    \item ¿Cómo puede utilizarse la visión por computadora para detectar obstáculos y objetos relevantes en el entorno de una persona con discapacidad visual?
    \item ¿Qué modelos de inteligencia artificial permiten realizar detección de objetos en tiempo real con suficiente precisión para aplicaciones de asistencia a la movilidad?
    \item ¿Cómo puede construirse un mapa del entorno y estimarse la posición del usuario dentro de él para proporcionar rutas e indicaciones auditivas sin depender de señal satelital?
    \item ¿De qué manera el reconocimiento de voz puede facilitar la interacción entre el usuario y el sistema sin depender de interfaces visuales?
    \item ¿Qué impacto tiene la implementación de un prototipo basado en inteligencia artificial en la seguridad y autonomía de las personas con visión reducida durante su desplazamiento?
\end{itemize}